\section{Introduction}
\label{sec:intro}

Think of a navigation service that suggests a route through the city. The travel time along each road changes throughout the day---congestion builds during rush hour, eases overnight, and varies from one day to the next. The service has access to historical speed data and can estimate what the roads will look like on average, but when the user actually drives, the realised travel times are a sample from that distribution. Which route should the service recommend?

A typical approach is to compute the shortest path with respect to the current (known) travel times, for instance by running Dijkstra's algorithm at query time. This returns the optimal path for the current conditions, every time, but it is expensive: each query costs $\mathcal O(m \log n)$ time, where $m$ is the number of edges and $n$ the number of vertices in the graph. If the navigation service handles many queries per second the computational load adds up. You would think that much of this computation is repetitive, even if slightly different every time. In principle the entire costs could be shifted to a preprocessing phase by computing and storing an answer for every possible traffic scenario, but the space required grows with the number of scenarios and quickly becomes impractical.

The problem of answering many shortest-path queries over a distribution of edge costs we refer to as the \emph{stochastic shortest path problem} (\cref{prob:shortest-path}). A recent line of work~\cite{miltos} defines the \emph{portfolio optimization} problem which our problem can be considered an instance of; during preprocessing we select a small set of candidate paths (a portfolio) for each source--target pair, so that in any traffic scenario at least one of them is cheap. At query time we only need to evaluate the portfolio under the current conditions and return the cheapest---an operation that is linear-time in the size of the portfolio and the size of its paths. The tradeoff between the size of the portfolio and the expected approximation quality is thus the central question.

This project explores this tradeoff on a real-world traffic dataset. We take the road network of a Shanghai district with travel speeds aggregated over 10-minute intervals throughout April 2015, which gives us a concrete distribution of edge cost functions. For this instance we implement three portfolio algorithms: a greedy algorithm that incrementally builds a portfolio by minimising expected cost over the full path set, a frequency-based selection that picks the paths that are shortest the most often, and a lightweight penalty scheme that repeatedly runs Dijkstra on expected edge costs while discouraging reuse of edges. We evaluate these algorithms in terms of the expected approximation factor they achieve for different portfolio sizes and the preprocessing and query times they require, and we compare them against the baseline of running Dijkstra at every query.

Our goal is not to claim that one algorithm is universally best, but rather to understand the practical behaviour of the portfolio model on a realistic input. We find that on the Shanghai dataset, portfolios of modest size ($k \approx 8$) achieve expected costs within hundredths of a percent of the optimal, and that the frequency-based approach---simple and fast---performs nearly as well as the greedy algorithm which has to iterate over the exponentially many candidate paths. The penalty-based scheme, while the cheapest to precompute, converges more slowly and illustrates the cost of ignoring the distributional information beyond the mean.
